\documentclass[20pt,extrafontsizes,article,oneside]{memoir}
%http://mirror.las.iastate.edu/tex-archive/macros/latex/contrib/memoir/memman.pdf
% see page 46 for font size options table 3.9
\usepackage[utf8]{inputenc}
\usepackage[ngerman]{babel}
\usepackage{amsmath}
\usepackage{amsfonts}
\usepackage{amssymb}
\usepackage{graphicx}
\usepackage{calc}
\usepackage{ccicons}
\usepackage{lipsum}

\usepackage{geometry}
\geometry{a0paper,landscape,left=17.5cm,right=4cm,top=2cm,bottom=2cm}
%\geometry{showframe=true} 
\usepackage{layout}

\usepackage{background}
\backgroundsetup{
	scale=1,
	color=black,
	opacity=1,
	angle=0,
	contents={%
		\includegraphics[width=\paperwidth,height=\paperheight]{parts/background-image}
	}%
}


\usepackage{chronosys}

\usepackage{xcolor}
\definecolor{faculty}{RGB}{0,61,143}

%change to iftrue to turn debug boxes on
\iffalse
	\newcommand{\debugBox}[1]{\fbox{#1}}
\else
	\newcommand{\debugBox}[1]{#1}
\fi

\author{Bennet Bleßmann, Henning Sachmann}
\title{}
\begin{document}


\textbf{\huge Was kann schief gehen bei der Entwicklung von algorithmischen Entscheidungsfindungsprozessen?}

\hfill

\debugBox{\begin{minipage}[t][0.3\textheight][t]{0.97\textwidth}
\Large 

	\debugBox{\begin{minipage}[t][0.3\textheight]{0.37\textwidth}
		\debugBox{\begin{minipage}[t][345pt]{430pt}
		\makebox[420pt]{}%empty minipage where the CAU-Logo bulges into the page
		\end{minipage}}
		\debugBox{\begin{minipage}[t][365pt]{545pt}
		%minipage right of the logo7
		\underline{Info}
% sadly lipsum befor version 2.0 does not have the option to generate less than a pragraph :(

		Im Rahmen dieses Posters wollen wir darstellen was bei der Entwicklung von Algorithmen zur Entscheidungsfindungsprozessen,
		im spezifischen beziehen wir uns hier auf Deep-Learning-Algorithmen, welche den für uns interresanten Teil darstellen.
		Für diesen Zweck Unterteilen wir die Entwicklung ähnlich dem Wasserfallmodel in einzelne Phasen hier Anforderund, Entwurf, Umsetzung, Training und 
		\end{minipage}}
		
		\debugBox{\begin{minipage}{975pt}
		
		%left third above the timeline under the bulge
		Anwendung. In jeder dieser Phasen können verschiedene Probleme auftreten. %TODO
		\end{minipage}}
	\end{minipage}}
	\debugBox{\begin{minipage}[t][0.3\textheight]{0.3\textwidth}
	\underline{Probleme im Bereich der Anforderungen}	
	
	\begin{itemize}
	\item 	
	\end{itemize}		
	\end{minipage}}
	\debugBox{\begin{minipage}[t][0.3\textheight]{0.3\textwidth}
	\underline{Right}	

	\lipsum[4-5]	
	\end{minipage}}

\end{minipage}}

%chronology
\begin{minipage}[t][0.3\textheight][t]{0.97\textwidth}
{\HUGE Ablauf bei der entwicklung einer Deep-Learning-Anwendung}
\startchronology[
	startyear=0,
	stopyear=16,
	dates=false, % don't display beginning and end date 
	height=100pt,
	color=faculty]

\setupchronoevent{date=false,markdepth=40pt,textstyle=\Large}
\setupchronoperiode{dates=false,bottomdepth=10pt,topheight=90pt,textdepth=50pt,textstyle=\Large}

\chronoperiode{0}{3}{Anforderung}
\chronoevent{1}{Kunden Anforderung}
\chronoevent[markdepth=-120pt]{2}{Machbarkeits prüfung}

\chronoperiode{3}{6}{Entwurf}

\chronoevent[markdepth=-120pt]{4}{Bestimmung der benötigten Daten}
\chronoevent{5}{Wahl/Entwurf des Algorithmus}

\chronoperiode{6}{8}{Entwicklung/Umsetzung des Entwurfs}

\chronoperiode[textdepth=-150pt]{8}{12}{Training}
\chronoevent{9}{Sammeln der Trainingsdaten}
\chronoevent{10}{Filtern der Trainingsdaten}
\chronoevent{11}{Training mit den Trainingsdaten}


\chronoperiode{12}{16}{Wartung/Anwendung}
\stopchronology
\end{minipage}

\debugBox{
\begin{minipage}[t][0.3\textheight-200pt][t]{0.97\textwidth}
\Large 
Content below time line
\end{minipage}
}


{
\color{white}
\huge

\llap{\makebox[16cm][c]{Autoren:}}

\llap{\makebox[16cm][c]{Bennet Bleßmann}}

\llap{\makebox[16cm][c]{Henning Sachmann}}
}

%TODO ask if this the following is ok if so change to iftrue or remove iffalse and fi
%creative commons with attribution and share-alike
\iftrue
{

\hfill

\large
{
\color{white}
\llap{\makebox[16cm]{\ccLogo \ccAttribution \ccShareAlike}}
}
Dieses Material steht unter der Creative-Commons-Lizenz Namensnennung - Weitergabe unter gleichen Bedingungen 4.0 International.

Um eine Kopie dieser Lizenz zu sehen, besuchen Sie http://creativecommons.org/licenses/by-sa/4.0/.
}
\fi
\clearpage

%un-comment to generate testpage
%\NoBgThispage
\layout
\clearpage


\end{document}